\documentclass{beamer}
\usetheme{metropolis}           % Use metropolis theme
\title{Food commodity prices, shocks and factor markets in Africa: Implications on income distribution, food and nutrition security}
\date{\today}
\logo{\includegraphics[height=0.8cm]{logo_zef.png}\vspace{220pt}}
\author{\includegraphics[height=1cm,width=5cm]{logo_zef.png} \\ Henry Kankwamba}
%\institute{Centre for Development Research (ZEF)}
\usepackage{Sweave}
\begin{document}
\Sconcordance{concordance:presentation.tex:presentation.Rnw:%
1 7 1 1 0 63 1}

%\SweaveOpts{concordance=TRUE}
  \maketitle
  \begin{frame}
\begin{itemize}
	\item     <1-| alert@1> Commodity price shocks could have significant impacts on factor markets and welfare (Abott and Batisti, 2009; Byerlee et al., 2006; Deaton and Miller, 1995). 
	\item <2-| alert@2> Higher prices could be beneficial to owners of factors of production by increasing their incomes (Breisinger et al., 2009). 
	\item <3-| alert@3> Higher income could stimulate demand for commodities and also lead to more savings and investment (Burfisher, 2011). Shocks e.g. price shocks adverse weather events could have reverse effects.
	\end{itemize}
  \end{frame}
  
  \begin{frame}
  	\begin{itemize}
  		\item <1-| alert@1> Of note, agricultural households are complex. There are net producers who maximize profit, net consumers who maximize utility and autarkic households (Sadoulet and de Janvry, 1995; Benjamin, 1992). 
  		\item A shock to these households could have varied results on these households. Producers could benefit, consumers would be hurt while autarkic households might be unaffected or might face indirect impacts.
	    \item  <2-| alert@2> Therefore, making judgements about welfare effects could be difficult in absence of data (Kalkuhl et al., 2016).  
  	\end{itemize}
  \end{frame}
  \begin{frame}
  	\begin{itemize}
  		\item <1-| alert@1> While there is evidence that markets are generally imperfect in Sub-Saharan Africa, little attention has been paid towards how commodity markets interact with factor markets (Rashid and Minot, 2010; Poulton et al., 2006; Barrett et al., 2001). 
  		\item It could be interesting to examine how household labor allocation decisions and shadow prices for labor change when shocks hit households.  
  		\item <2-| alert@2> Further, how policy and distortions in agricultural commodity markets could influence factor market or vice versa has also not been addressed adequately. 
  	\end{itemize}
  \end{frame}
  \begin{frame}
  	\begin{itemize}
  		\item <1-| alert@1> Most studies focus on commodity market interactions (Minot et al., 2011; Norton, 2004; Jayne et al., 2002; Readon et al., 1994). 
  		\item <2-| alert@2>Further, effects of market interactions on income distribution, food and nutrition security outcomes among heterogeneous households has not been thoroughly examined. Notable examples usually use descriptive measures (Case, 1996; Barret and Dorosh, 1998).
  	\end{itemize}
  \end{frame}
  
  \begin{frame}
  The main objective of this research is to assess the impact of staple food price shocks on factor markets, household income distribution, food and nutrition security. 
  
Specifically, it seeks 
  	\begin{enumerate}
  		\item <1-| alert@1> To quantify effects of price policy and exogenous distortions on local staple food markets. 
  		\item <2-| alert@2> To assess the impact of a staple food price shock on labor allocation decisions and shadow prices of labor.
  		\item <3-| alert@3>To measure the economy wide distributional effects of a staple food price shock on real consumption, calorie and micronutrient intake.  
  	\end{enumerate}
  
  \end{frame}
  
  \begin{frame}
  The study seeks to answer  these research questions.
  \begin{itemize}
  \item <1-| alert@1> To what extent do staple food price policy and exogenous distortions reflect on local food markets?
  \item <2-| alert@2> How much do staple food price changes affect the effective wage in the medium term?
\item	<3-| alert@3> Can heterogeneity in staple food preferences reflect vulnerability to food price shocks? 
\item  <4-| alert@4> To what extent do shocks affect real consumption, calorie and micronutrient intake?  
  \end{itemize}
  \end{frame}
  
\begin{frame}{Methods}
In general, I am thinking in an agricultural household modeling framework
\begin{itemize}
\item To achieve objective one, the study will use threshold regression approaches.
\item To achieve objective two, the study will use sample selection regression models.
\item To achieve the third objective, the study will use a demand systems approach and microsimulations.
\end{itemize}
\end{frame}
\end{document}}
